\section{Preliminaries}
\subsection*{Fuzzy Image Processing}
Fuzzy image processing encompasses a range of techniques designed to understand, represent, and process images, along with their segments and features, as {\em fuzzy sets} (see Figure \ref{fig:fii}). The process begins with the {\em fuzzification} (or encoding) of image data, followed by the {\em defuzzification} (or decoding) of the results. These steps enable the application of fuzzy techniques to image processing. At the core of fuzzy image processing lies {\em membership modification}~\cite{haussecker2000fuzzy}. The workflow of fuzzy image processing is illustrated in~\Cref{fig:fii}.

\begin{figure*}
    \centering
    \begin{tikzpicture}[
        node distance=2cm and 3cm,
        every node/.style={draw, align=center, rounded corners},
        every edge/.style={draw, -{Latex[length=3mm]}}
    ]
    
    % Nodes
    \node (expert) {expert knowledge};
    \node (fuzzification) [below left=of expert] {image\\ fuzzification};
    \node (modification) [below=of expert] {membership\\ modification};
    \node (defuzzification) [below right=of expert] {image\\ defuzzification};
    \node (fuzzy) [below=of modification] {fuzzy logic \\ fuzzy set theory};
    
    % Edges
    \draw[->] (expert) -- (fuzzification);
    \draw[->] (expert) -- (modification);
    \draw[->] (expert) -- (defuzzification);
    \draw[->] (fuzzification) -- (modification);
    \draw[->] (modification) -- (defuzzification);
    \draw[->] (fuzzy) -- (fuzzification);
    \draw[->] (fuzzy) -- (modification);
    \draw[->] (fuzzy) -- (defuzzification);
    \end{tikzpicture}
    \caption{The workflow of fuzzy image processing}
    \label{fig:fii}
\end{figure*}
\subsection*{Fuzzy Sets for Intensity Transformation}
Contrast enhancement is one of the principal
applications of intensity transformations. We can state the process of enhancing the contrast of a gray-scale image using the following rules \cite{gonzalez2008digital}:
\begin{itemize}
    \item IF a pixel is \textbf{dark}, THEN make it \textbf{darker}
    \item IF a pixel is \textbf{gray}, THEN make it \textbf{gray}
    \item IF a pixel is \textbf{bright}, THEN make it \textbf{brighter}
\end{itemize}